\chapter*{Abstract}
\addcontentsline{toc}{chapter}{Abstract}

Modern video sharing platforms must solve a common core problem: when a user uploads a
source file in an arbitrary format, codec, or resolution, the system has to transcode it
automatically into HTTP Live Streaming (HLS) so that browsers can play it with adaptive
quality. Provisioning a dedicated server cluster or a Kubernetes system that runs
continuously purely to wait for transcoding work is wasteful, because upload traffic is
inherently bursty and the infrastructure sits idle most of the time. Conversely, when many
users upload simultaneously, the system must scale out to process hundreds of videos in
parallel.

This project presents the design, implementation, and evaluation of a complete video
sharing platform whose transcoding subsystem is built on a Serverless Container
and Event-Driven architecture. Uploading a file to Amazon S3 emits an event that
is buffered in Amazon SQS and dispatched by an AWS Lambda job submitter to AWS Batch
running on AWS Fargate. A container packaging FFmpeg is created only when there is work to
do, produces three renditions at 360p, 720p, and 1080p together with a master playlist and
a thumbnail, uploads the result to a processed bucket, updates the metadata in MongoDB
Atlas, and then terminates. Content is delivered through Amazon CloudFront with Origin
Access Control, and every stream is gated at the edge by CloudFront Signed Cookies, so that a
viewer without a valid signature is rejected even when they know the exact manifest URL; a
video's visibility setting governs who the backend will issue those cookies to, not whether
the gate applies.

The entire AWS infrastructure is described as code using eleven independent Terraform
modules and can be provisioned with a single command. Seven GitHub Actions workflows
automate linting, unit testing, static security analysis, container image scanning,
infrastructure validation, and deployment, with a two-tier quality gate that blocks a pull
request whenever a fixable critical vulnerability is detected.

The system is evaluated along four dimensions: quality of experience measured through
Time-to-First-Frame, transcoding pipeline latency across 100~MB, 500~MB, and 1~GB inputs,
scalability under a concurrent upload stress test, and operating cost compared against a
traditional always-on Amazon EC2 deployment. The cost analysis shows that, for the
architecture-dependent portion of the bill, the serverless approach removes almost all idle
resource waste, although this advantage narrows as sustained utilisation increases.
